\section{Semilinear spread rigidity and the one-dimensional branch}
The following full argument is reused from the reviewed phase4 theorem
\emph{Rigidity of locally primitive semilinear spreads}. Its application
to the source question is already recorded there; this edition claims
no new question answer or priority. We retain Delandtsheer's reduction
quoted as source Proposition 5.2: a one-dimensional flag-transitive action
contains all translations, and its spread case consists of cosets of
equal-dimensional additive subspaces. The proof below includes the
restriction kernels and characteristic-dividing Galois case.

\begin{lemma}[Prime degree and arithmetic]\label{prime}
A finite primitive permutation group with cyclic normal subgroup and
cyclic quotient has prime degree, if its degree exceeds one. If
\[
 r=\frac{p^{u\ell}-1}{p^\ell-1}
\]
is prime, where $p$ is prime and $u>1$, then $u$ is prime,
$\ell=u^e$ for some $e\ge0$, and $p\not\equiv1\pmod u$.
\end{lemma}
\begin{proof}
Pass to the faithful induced action. The cyclic normal subgroup either
is trivial or is transitive. In the former case the whole induced group
is cyclic, hence regular of prime order. In the latter case the cyclic
normal subgroup is regular. Any proper nontrivial characteristic subgroup
would give a nontrivial block system preserved by the whole group, so
its order and the degree are prime.

If $u$ is composite, its geometric sum factors into two integer
geometric sums greater than one. Thus $u$ is prime. If a prime
$v\ne u$ divides $\ell$, put $x=p^{\ell/v}$. Reduction modulo
$\Phi_u(X)=1+X+\cdots+X^{u-1}$ shows that $\Phi_u(X)$ divides
$\Phi_u(X^v)$: multiplication by $v$ permutes the nonzero residues
modulo $u$. It follows that
$1<\Phi_u(x)<\Phi_u(x^v)=r$, contradicting primality. Hence
$\ell=u^e$. Finally $p\equiv1\pmod u$ would give $u\mid r$ while
$r>u$, again impossible.
\end{proof}

\begin{lemma}[Cyclic semilinear descent]\label{descent}
Let $p,u$ be primes with $p\not\equiv1\pmod u$. Let a finite group $J$
act faithfully and irreducibly on an $\F_p$-space $U$ of dimension $u^e$.
Suppose $C\triangleleft J$ is cyclic, $K=\F_p[C]$ is a field of
endomorphisms, $J$ normalizes $K$, and $J/C$ is a cyclic $u$-group.
Then $\dim_KU=1$.
\end{lemma}
\begin{proof}
Write $|K|=p^h$. Since $h\mid u^e$, $h$ is a power of $u$.
For every power $u^a$, the number $p^{u^a}-1$ is prime to $u$:
this is immediate for $p=u$, and otherwise follows from
$p^{u^a}\equiv p\not\equiv1\pmod u$. In particular $C$ has order
prime to $u$.

Choose a lift $t$ of a generator of $J/C$ and let $y$ be its
$u$-primary component inside the cyclic group generated by $t$.
The prime-to-$u$ component lies in $C$, so $J=C\langle y\rangle$.
Let $\sigma$ be the automorphism of $K$ induced by $y$ and let $m$
be its order. Then $z=y^m$ centralizes $K$, $C$ and $y$, hence $J$.
Every nonzero element of the commutative algebra $\F_p[z]$ is invertible:
its kernel is $J$-invariant, so irreducibility makes it injective. This
finite algebra is a field of order $p^a$, with $a\mid u^e$. Its
multiplicative group has order prime to $u$. Since $z$ has $u$-power
order, $z=1$. Thus the order of $y$ is exactly $m$.

There is a nonzero $y$-fixed vector without dividing by $m$. Choose
$0\ne v\in U$ and consider
\[
 S(b)=\sum_{i=0}^{m-1}\sigma^i(b)y^i(v),\qquad b\in K.
\]
Every $S(b)$ is fixed by $y$, by cyclically shifting the summands.
They cannot all vanish. Apply a $K$-linear functional nonzero on $v$;
vanishing would give a nonzero linear relation over $K$ among the
distinct field automorphisms $1,\sigma,\ldots,\sigma^{m-1}$.
Their linear independence follows by taking a shortest relation and
subtracting its value at $bx$ from a suitable multiple of its value
at $x$, shortening the relation. For some $b$ we therefore obtain a
nonzero fixed vector $w=S(b)$. Its $K$-line is $J$-invariant: $C$
acts by scalars and $y$ normalizes $K$ while fixing $w$.
Irreducibility gives $U=Kw$.
\end{proof}

\begin{theorem}[Locally primitive one-dimensional spread spaces]\label{spread}
Let $F=\F_{p^d}$, let $\mathcal S$ be a spread of $\ell$-dimensional
additive subspaces, with $0<\ell<d$, and let
$H\le\Gamma L_1(F)$ preserve it and act transitively. Suppose $H$
is primitive on $\mathcal S$ and $H_U$ is irreducible on a spread
member $U$. Then $\mathcal S$ is exactly the set of one-dimensional
$\F_{p^\ell}$-subspaces of $F$. Its size is prime and
$d=u^{e+1},\ell=u^e$ for a prime $u$.
Conversely such a locally primitive one-dimensional action exists
precisely when $(p^d-1)/(p^\ell-1)$ is prime; the geometry is then
uniquely $\operatorname{AG}(d/\ell,p^\ell)$ up to isomorphism.
This does not classify every possible acting subgroup.
\end{theorem}
\begin{proof}
The normal scalar subgroup of $H$ is cyclic and its quotient embeds
in the cyclic Galois group. Lemma~\ref{prime} gives prime spread size
$r=(p^d-1)/(p^\ell-1)$. Divisibility of the denominator into the
numerator implies $\ell\mid d$: dividing $d=a\ell+b$, $0\le b<\ell$,
would otherwise give $p^\ell-1\mid p^b-1$. Apply the arithmetic part
of that lemma with $u=d/\ell$.

Put $C_0=H_U\cap F^*$, and let $J$ be the faithful image of $H_U$
on $U$. The ambient scalar group $C_0$ restricts faithfully on nonzero
$U$, and $K=\F_p[C_0]$ is an actual subfield of $F$. Also $U$ is a
$K$-space. The semilinear stabilizer normalizes $K$; its faithful
image modulo the image $C$ of $C_0$ is a quotient of a subgroup of
$\operatorname{Gal}(F/\F_p)$, hence a cyclic $u$-group because
$d=u^{e+1}$. Lemma~\ref{descent} now gives $U=wK$ with
$|K|=p^\ell$. The unique subfield of that order is $\F_{p^\ell}$.
All $H$-images of $U$ are $K$-lines, since scalar-Frobenius maps
preserve this subfield. There are exactly $r$ such lines, so these
are the entire spread.

For the converse use $F^+$\,$\rtimes F^*$ on the $K$-line spread.
It is flag-transitive, its pencil action is the cyclic regular
prime-degree action of $F^*/K^*$, and its line stabilizer contains
$K^+\rtimes K^*$, sharply two-transitive on $K$. Its global action
on $F$ is sharply two-transitive as well. The local-action criterion
proves every required property. This establishes existence within a
one-dimensional group, without borrowing the full automorphism group
of the affine space.
\end{proof}

\begin{proposition}[Local two-transitivity in the one-dimensional branch]\label{onedim2t}
Under the spread hypotheses, local two-transitivity occurs only for
$\operatorname{AG}(2,2)$ with $A\Gamma L_1(4)$, or
$\operatorname{AG}(2,4)$ with $A\Gamma L_1(16)$.
\end{proposition}
\begin{proof}
Write $v=p^{u^{e+1}}$, $k=p^{u^e}$ and $d=u^{e+1}$ by
Theorem~\ref{spread}. Pencil two-transitivity forces
$r(r-1)\mid d(v-1)$. Since $v-1=r(k-1)$, we have
$r-1\mid d(k-1)$. But $k\mid r-1$ and $\gcd(k,k-1)=1$,
so $k\mid d$. Hence $p=u$ and $u^e\le e+1$, leaving only $e=0$
or $e=1$; in the latter case $p=u=2$. If $e=0$, cancellation gives
$1+p+\cdots+p^{p-2}\mid p-1$, which is possible only for $p=2$.
Thus the only parameter pairs are $(v,k)=(4,2),(16,4)$.

For four points, pencil two-transitivity requires a group order
divisible by $4\cdot3\cdot2=24$, so the full $A\Gamma L_1(4)$ is
necessary and is $S_4$. For sixteen points, the five-point pencil
requires $16\cdot5\cdot4\mid |G|$, and the twenty lines with four
points each require $20\cdot4\cdot3\mid |G|$. Their least common
multiple is $960$, exactly $|A\Gamma L_1(16)|$; the full group is
necessary. It suffices: on the five cosets $\F_{16}^*/\F_4^*$ the
multipliers are a regular cyclic group of order five and Frobenius
acts by multiplication by two, of order four on that quotient.
This is sharply two-transitive. On a line $\F_4$, its translations
and nonzero scalar multipliers give a sharply two-transitive group.
\end{proof}
